\chapter{共形映射 (Conformal Mapping)}

\section{几何意义}
\begin{itemize}
    \item \textbf{保角性}：旋转角 $\Arg f'(z_0)$ 描述了曲线经过映射后的旋转量，且具有\textbf{旋转角不变性}。
    \item \textbf{伸缩率不变性}：$|f'(z_0)|$ 描述了点 $z_0$ 处的微小线段在映射后的伸缩比例。
\end{itemize}

\section{共形映射的定义}
若 $w = f(z)$ 在 $z_0$ 处解析、一一对应且 $f'(z_0) \neq 0$，则称该映射在 $z_0$ 处是\textbf{共形的}。

\subsection{扩充复平面上的共形性}
针对扩充复平面 $\mathbb{C}^*$ 上的定义：
\begin{itemize}
    \item \textbf{情形 I}：当 $w = f(z)$ 将 $z = z_0 (\neq \infty)$ 映成 $w = \infty$ 时：
    作 $\zeta = \dfrac{1}{w}$。若 $\zeta = \dfrac{1}{f(z)}$ 在 $z = z_0$ 处是共形的，则称 $w = f(z)$ 在 $z = z_0$ 处是共形的。
    \item \textbf{情形 II}：当 $w = f(z)$ 将 $z = \infty$ 映成 $w = w_0 (\neq \infty)$ 时：
    作 $\eta = \dfrac{1}{z}$。若 $w = f\left(\dfrac{1}{\eta}\right)$ 在 $\eta = 0$ 处是共形的，则称 $w = f(z)$ 在 $z = \infty$ 处是共形的。
    \item \textbf{情形 III}：当 $w = f(z)$ 将 $z = \infty$ 映成 $w = \infty$ 时：
    作 $\eta = \dfrac{1}{z}$，$\zeta = \dfrac{1}{w}$。若 $\zeta = \dfrac{1}{f(1/\eta)}$ 在 $\eta = 0$ 处是共形的，则称 $w = f(z)$ 在 $z = \infty$ 处是共形的。
\end{itemize}

\begin{theorem}[重要]
分式线性映射在扩充复平面上是一一的，且具有\textbf{保角性}。
\end{theorem}

\section{分式线性映射}
分式线性映射的一般形式为
\[
w = \frac{az+b}{cz+d}, \quad ad - bc \neq 0.
\]
它可以分解为以下三种基本变换：
\begin{enumerate}
    \item \textbf{平移}：$z \mapsto z + b$
    \item \textbf{旋转与伸缩}：$z \mapsto az$
    \item \textbf{反演}：$z \mapsto \dfrac{1}{z}$
\end{enumerate}
公式分解：
\[
w = \frac{a}{c} - \frac{ad-bc}{c} \cdot \frac{1}{cz+d}.
\]

\section{保圆性与对称点}
\subsection{对称点引理}
\begin{lemma}
$z_1, z_2$ 是关于圆周 $C$ 的对称点的充要条件是通过 $z_1, z_2$ 的任何圆周 $L$ 与 $C$ 正交。
\end{lemma}

\textbf{证明思路：}
\begin{itemize}
    \item 设 $C$ 的圆心为 $z_0$，半径为 $R$。
    \item 对称点满足：$|z_1 - z_0| \cdot |z_2 - z_0| = R^2$。
    \item 利用切割线定理可证正交性。
\end{itemize}

\subsection{保圆性 (Circle-Preserving Property)}
\begin{definition}
若一个映射将圆周映射为圆周，则称它具有\textbf{保圆性}。
\end{definition}

以反演映射 $w = \dfrac{1}{z}$ 为例：
令 $z = x + iy$，$w = u + iv$，代入一般圆方程
\[
a(x^2+y^2) + bx + cy + d = 0,
\]
可得
\[
d(u^2+v^2) + bu - cv + a = 0,
\]
仍为圆方程，故保圆。

\begin{remark}
分式线性映射在复平面上将圆周映成圆周，具有\textbf{保圆性}。
\end{remark}

\section{交比与映射唯一性}
\subsection{交比 (Cross Ratio)}
交比的定义为
\[
(z_1, z_2, z_3, z_4) = \frac{z_1 - z_3}{z_2 - z_3} : \frac{z_1 - z_4}{z_2 - z_4}
= \frac{(z_1 - z_3)(z_2 - z_4)}{(z_2 - z_3)(z_1 - z_4)}.
\]

\begin{remark}
分式线性映射具有\textbf{保交比性}。
\end{remark}

\subsection{唯一确定映射的条件}
\begin{theorem}
在 $z$ 平面与 $w$ 平面各给定相异三点，则存在唯一的分式线性映射 $w = f(z)$ 满足三点对应关系。
\end{theorem}

计算公式：
\[
\frac{(w - w_1)(w_3 - w_2)}{(w - w_2)(w_3 - w_1)}
= \frac{(z - z_1)(z_3 - z_2)}{(z - z_2)(z_3 - z_1)}.
\]

\section{三种典型映射类型及其推导}
分式线性映射的核心在于利用\textbf{边界对应原理}与\textbf{对称点原理}确定参数。

\subsection{类型 I：上半平面 $\im(z) > 0 \to$ 单位圆盘 $|w| < 1$}
\textbf{推导：}
\begin{enumerate}
    \item 设上半平面内一点 $\lambda$（$\im(\lambda) > 0$）映为圆心 $w=0$。
    \item 由对称点原理，$\lambda$ 关于实轴的对称点 $\bar{\lambda}$ 应映为 $w=\infty$。
    \item 映射形式为 $\displaystyle w = k \,\frac{z - \lambda}{z - \bar{\lambda}}$。
    \item 实轴映为单位圆周 $|w|=1$，故 $|k|=1$，即 $k = e^{i\theta}$。
\end{enumerate}

\[
w = e^{i\theta} \frac{z - \lambda}{z - \bar{\lambda}}, \quad \theta\in\mathbb{R},\ \im(\lambda) > 0.
\]

\subsection{类型 II：单位圆盘 $|z| < 1 \to$ 单位圆盘 $|w| < 1$}
\textbf{推导：}
\begin{enumerate}
    \item 设圆盘内一点 $\alpha$（$|\alpha| < 1$）映为 $w=0$。
    \item 对称点 $1/\bar{\alpha}$ 映为 $w=\infty$。
    \item 映射形式为 $\displaystyle w = k' \,\frac{z - \alpha}{1 - \bar{\alpha}z}$。
    \item 边界对应得 $|k'|=1$，即 $k' = e^{i\varphi}$。
\end{enumerate}

\[
w = e^{i\varphi} \frac{z - \alpha}{1 - \bar{\alpha}z}, \quad \varphi\in\mathbb{R},\ |\alpha| < 1.
\]

\subsection{类型 III：上半平面 $\im(z) > 0 \to$ 上半平面 $\im(w) > 0$}
\textbf{推导：}
\begin{enumerate}
    \item 系数 $a,b,c,d$ 均取实数。
    \item 虚部满足
    \[
    \im(w) = \frac{1}{2i}(w - \bar{w}) = \frac{(ad-bc)\im(z)}{|cz+d|^2}.
    \]
    \item 定向要求 $ad - bc > 0$。
\end{enumerate}

\[
w = \frac{az + b}{cz + d}, \quad a,b,c,d\in\mathbb{R},\ ad-bc > 0.
\]

\section{初等函数的共形性}
\subsection{幂函数}
幂函数 $w = z^n \ (n \ge 2)$：
\begin{itemize}
    \item 导数：$\dfrac{\dif w}{\dif z} = nz^{n-1}$。
    \item 共形性：$z \neq 0$ 时处处共形；$z=0$ 处角度放大 $n$ 倍。
    \item 角形域变换：将 $0 < \arg z < \alpha$ 映为 $0 < \arg w < n\alpha$。
\end{itemize}

\subsection{指数函数}
指数函数 $w=e^z$：
\begin{itemize}
    \item 导数：$\dfrac{\dif w}{\dif z}=e^z\neq0$，复平面内处处共形。
    \item 几何意义：将水平带形域 $0<\im z < 2\pi$ 共形映射为沿正实轴割开的复平面。
\end{itemize}